\section{The well-specified logistic model}\label{sec:well-sepcified-mod}

We begin in a setting that allows the cleanest and most precise comparisons
between a method using aggregated labels~\eqref{eq:logistic-majority} and
one without~\eqref{eq:maximum-likelihood} by considering the logistic link
for~\eqref{eq:general-calibration-model},
\begin{align*}
  \sigma\lr(t) = \frac{1}{1 + e^{-t}} \in \fsymlink.
\end{align*}
%This simplicity allows results that highlight many of the conclusions we
%draw, and so we present initial, relatively clean, results for the
%estimators~\eqref{eq:maximum-likelihood} and~\eqref{eq:logistic-majority}
%here.  
In particular, we assume identical links $\sigma_1\opt = \dots =
\sigma_m\opt = \sigma\lr$ and an i.i.d.\ sample $X_1, \ldots, X_n$, where
for each $i \in [n]$ we draw $(Y_{i1}, \ldots, Y_{im}) \simiid
\Prb_{\sigma\lr,\theta\opt}(\cdot \mid X_i)$ for a true vector $\theta\opt$.


\subsection{The isotropic Gaussian case} \label{sec:gauss}

To deliver the general taste of our results on the performance of the full
information~\eqref{eq:maximum-likelihood} and majority
vote~\eqref{eq:logistic-majority} approaches, we start by studying the
simplest case when $X \sim \normal(0, I_d)$.

\paragraph{Performance with non-aggregated data.}

The
non-aggregated MLE $\tmle$ in Eq.~\eqref{eq:maximum-likelihood}
admits a standard analysis,
which we state as a corollary of
Proposition~\ref{prop:lowd-maximum-likelihood} to come.
%
%Let the population log loss be $L(\theta) = \Prb \loss_\theta(Y \mid X)$. When $m$ is fixed, we can view $(X_i, Y_{i1}, \cdots, Y_{im})$ as independent samples from $\R^d \times \left\{\pm 1\right\}^m$ which allows us to directly employ the standard asymptotic normality for maximum likelihood estimators. Let $\what{u}\lr_{n,m}$ be the normalized version of $\what{\theta}\lr_{n,m}$. The classical theory tells us the following.
\begin{corollary}
  \label{cor:lowd-maximum-likelihood-Gaussian} 
  Let $X \sim \normal(0, I_d)$ and $t\opt = \ltwos{\theta\opt}$. The maximum
  likelihood estimator $\tmle$ is consistent, with $\tmle \cp
  \theta\opt$, and for $\proj_{u\opt}^\perp = I_d - u\opt {u\opt}^\top$,
  \begin{align*}
    \sqrt{n} (\what{u}\lr_{n,m} - u^\star)
    & \cd \normal\left(0,  m^{-1}
    \cdot (t\opt)^{-2} \cdot \E[\sigma\lr(\<\theta\opt, X\>)
      (1 - \sigma\lr(\<\theta\opt, X\>))]^{-1}
    \proj_{u\opt}^\perp \right).
  \end{align*}
\end{corollary}
\noindent
So the non-aggregated MLE is calibrated, recovering 
$\theta\opt$,
and enjoys convergence rates scaling as
$1/\sqrt{nm}$, so that increasing the number of labels $m$ linearly 
improves convergence rate.


\paragraph{Performance with majority-vote aggregation.}

The analysis of the majority vote estimator~\eqref{eq:logistic-majority}
requires more care, though the assumption that $X \sim \normal(0, I_d)$ allows
us to calculate the limits explicitly. In
brief, we show that when $X$ is Gaussian, the estimator is mis-calibrated
and its classification error converges more slowly in $m$
than the
non-aggregated classifier.
The basic idea is that the majority vote classifier
should still point in the direction of $\theta\opt$,
but it should be  ``calibrated'' to the probability of a majority
of $m$ labels being correct.

We follow this idea and sketch the derivation here, as it is central to all
of our theorems, then state a companion corollary.
%% to Corollary~\ref{cor:lowd-maximum-likelihood-Gaussian}.  
Each result depends on the probability
\begin{equation*}
  \Prb(\majY = \sgn(\<x, \theta^\star\>) \mid x)
  = \rho_m(|\<x, \theta^\star\>|)
\end{equation*}
of obtaining a correct label using majority vote,
where $\rho_m$ is the binomial probability
\begin{align}
  \rho_m(t) =
  \P\left(\bindist\Big(m, \frac{1}{1 + e^{-|t|}}\Big)\ge \frac{m}{2}\right)
  = 
  \sum_{i=\lceil m/2 \rceil}^m \binom{m}{i} \left(\frac{1}{1 + e^{-|t|}} \right)^i \left(\frac{e^{-|t|}}{1 + e^{-|t|}} \right)^{m-i}, \label{eq:link-function-m-majority-vote}
\end{align}
when $m$ is odd. (When $m$ is even the sum has the
asymptotically negligible additional additive
term $\frac{1}{2} \binom{m}{m/2} \frac{e^{-m|t|/2}}{(1 + e^{-|t|})^m}$.)
Key to the coming result is a
parameter that approximately equalizes binomial (majority vote) and logistic
(Bernoulli) probabilities, so we define
\begin{align}
  h_m(t) &  = \Ep \brk{|Z|(1-\rho_m(t^\star |Z|))}-\Ep
  \brk{\frac{|Z|}{1 + e^{t|Z|}}}, \quad \text{for }Z \sim \normal(0, 1).
  \label{eqn:h-centering}
\end{align}
We use $h$ to find the minimizer of population loss $\poploss\mv_m(\theta) =
\Ep[ \ell_\theta\lr(\majY \mid X)]$ via the ansatz
$\theta = t u^\star$ for
some $t > 0$.
%
Using the definition~\eqref{eq:link-function-m-majority-vote}
of $\rho_m$, for
$S\subopt = \sgn(\<X, \theta\opt\>)$,
\begin{align*}
  \poploss\mv_m(\theta)  
  & = \Ep\left[\log(1 + e^{-S\subopt \<X, \theta\>})
    \cdot\rho_m(t^\star |\<X, u^\star\>|) \right]
  + \Ep\left[\log(1 + e^{S\subopt\<X, \theta\>})
    \cdot (1-\rho_m(t^\star|\<X, u^\star\>|)) \right],
\end{align*}
where we recall $t\opt = \ltwo{\theta\opt}$, and compute the gradient
for $\theta = t u\opt$ via
\begin{align*}
  \nabla \poploss\mv_m(\theta) & =
  - \Ep \left[\frac{S\subopt \rho_m(t^\star |\<X, u^\star\>|) }{1 + \exp(S\subopt \<X, \theta\>)} X \right]
  + \Ep \left[\frac{S\subopt \exp(S\subopt \<X, \theta\>)
      (1-\rho_m(t^\star|\<X, u^\star\>|))}{1 + \exp(S\subopt \<X, \theta\>)}  X \right].
\end{align*}
Setting $Z = \<X, u^\star\>$ to decompose $X$ into the independent sum $X =
(X - u^\star Z) + u^\star Z$,
\begin{align}
  \nabla \poploss\mv_m(t u^\star) 
  & \stackrel{\mathrm{(i)}}{=} - \Ep \left[\frac{\sgn(Z)}{1 + \exp(t|Z|)} \rho_m(t^\star |Z|) X \right]  + \Ep \left[\frac{\sgn(Z)\exp(t|Z|)}{1 + \exp(t|Z|)} (1-\rho_m(t^\star |Z|)) X \right] \nonumber	\\
  %&= - \Ep \left[\frac{\sgn(Z)Z}{1 + \exp(t|Z|)} \rho_m(t^\star |Z|)\right] u^\star  + \Ep \left[\frac{\sgn(Z)Z\exp(t|Z|)}{1 + \exp(t|Z|)} (1-\rho_m(t^\star |Z|))\right]  u^\star \nonumber \\
  %& = \left(\Ep \left[\frac{|Z|  \exp(t|Z|)}{1 + \exp(t|Z|)}\right] - \Ep \left[|Z| \rho_m(t^\star |Z|) \right]\right) u^\star \nonumber \\
  & = \left(\Ep \brk{|Z|(1-\rho_m(t^\star  |Z|))}-\Ep \brk{\frac{|Z|}{1 + e^{t|Z|}}}\right) u^\star = h_m(t) u^\star  , \label{eq:majority-vote-root}
\end{align} 
where in (i) we substitute $S\subopt = \sgn(Z)$.
%
As we present in Corollary~\ref{cor:lowd-majority-vote-m}, $h_m(t)=0$ has a
unique solution $t_m \asymp \sqrt m$, so $t_m u\opt$ necessarily uniquely
minimizes the population loss $\poploss\mv_m$.

By completing the calculations for the value of $t_m$ above and leveraging
asymptotic normality, we have the following
special case of Proposition~\ref{prop:lowd-majority-vote-m-noise} to come.

\begin{corollary} \label{cor:lowd-majority-vote-m}
  Let $X \sim \normal(0, I_d)$
  and $t\opt = \ltwos{\theta\opt}$.
  There are numerical constants $a, b > 0$ such that for the function $h = h_m$ in~\eqref{eqn:h-centering},
  there is a unique $t_m \geq t_1 = t^\star$ solving $h(t_m) = 0$
  and 
  \begin{equation*}
    \what{\theta}\mv_{n,m} \cp t_m u\opt
    ~~ \mbox{and} ~~
    \lim_{m \to \infty} \frac{t_m}{t\opt \sqrt{m}} = a.
  \end{equation*}
  Moreover,
  there exists
  $C_m(t) = \frac{b}{t \sqrt{m}}(1 + o_m(1))$
  such that
  $\what{u}\mv_{n,m} = \mve/\ltwos{\mve}$ satisfies
  \begin{align*}
    \sqrt{n} \left(\what{u}\mv_{n,m} - u^\star \right) &\cd \normal\left(0, C_m(t\opt) {\proj_{u\opt}^\perp}\right).
  \end{align*}
\end{corollary}

It is instructive to compare the rates of this estimator to the rates
for the non-aggregated MLE
in \Cref{cor:lowd-maximum-likelihood-Gaussian}. First, the non-aggregated
estimator is calibrated in that $\tmle \to \theta\opt$, in contrast to the
majority-vote estimator, which roughly
``calibrates'' to the probability majority vote is correct
(cf.~\eqref{eq:majority-vote-root}) via the convergence
$\mve \to c \sqrt{m} \theta\opt$ as $n
\to \infty$.
%
The scaling of $C_m$ in
\Cref{cor:lowd-majority-vote-m} also shows that the majority-vote
estimator exhibits worse convergence rates by a factor of $\sqrt{m}$ than
the estimator $\what{\theta}\lr_{n,m}$: regardless of the dimension $d$,
for constants $c\lr$ and $c\mv$ 
that depend only on $t\opt = \ltwos{\theta\opt}$,
the asymptotic variances differ by $\sqrt{m}$, as
\begin{equation*}
  \sqrt{n}(\what{u}\lr_{n,m} - u\opt) \cd
  \normal\left(0, m^{-1} c\lr \proj_{u\opt}^\perp  \right)
  ~~ \mbox{while} ~~
  \sqrt{n}(\what{u}\mv_{n,m} - u\opt) \cd
  \normal\left(0, m^{-1/2} c\mv \proj_{u\opt}^\perp \right).
\end{equation*}

To obtain intuition for this result via comparison with
Corollary~\ref{cor:lowd-maximum-likelihood-Gaussian}, consider the Fisher
information $\E[\sigma\lr(\<\theta\opt, X\>) (1 - \sigma\lr(\<\theta\opt,
  X\>)) XX^\top]$ for logistic regression. Letting $Z\opt = \<\theta\opt,
X\>$ be the predicted margin, there is only ``information'' available for an
estimator when $\sigma\lr(Z\opt)$ is near $\half$, as otherwise
$\sigma\lr(Z\opt)(1 - \sigma\lr(Z\opt)) \approx 0$; under majority vote, we
analogously require $\P(\bindist(m, \sigma\lr(Z\opt)) > \frac{m}{2}) \approx
\half$, which in turn means that we need $\sigma\lr(Z\opt) \in \half \pm
O(\frac{1}{\sqrt{m}})$ by standard binomial concentration, or $Z\opt =
\<\theta\opt, X\> \in \pm O(\frac{1}{\sqrt{m}})$.  This occurs on about
$1/\sqrt{m}$ fraction of the sample, so of $n \cdot m$ total observations,
majority vote ``loses'' a factor $\sqrt{m}$.

\newcommand{\mt}{$\textsc{M}_\beta$}

\subsection{Comparisons for more general feature distributions}
\label{sec:non-gauss}

The key to the preceding results
is that $X$ decomposes into components aligned with
$\theta\opt$ and independent noise.
%
By abstracting away the Gaussianity assumptions, we may carefully track
margins and label noise, and their roles in the merits of
maximum-likelihood-type (full label information) estimators relative to
those using majority-vote labels.
%
The results follow from the master theorems in the sequel, which rely on the
following independence
\begin{assumption}
  \label{assumption:x-decomposition}
  The covariates $X$ have non-singular covariance $\Sigma$ and decompose as a
  sum of independent random vectors in the span of $u\opt$ and its
  complement
  \begin{equation*}
    X = W +  Z u^\star  ,
    ~~~\mbox{where}~~~
    W \perp\!\!\!\!\perp Z  , ~~
    \<W,u\opt\> = 0  ,
    ~ \Ep[W] = 0  ,
    ~ \Ep[Z] = 0  .
    %\label{eqn:X-distribution}
  \end{equation*}
\end{assumption}
%Under Assumptions~\ref{assumption:x-decomposition}, we characterize the limiting
%behavior of the majority vote and non-aggregated models,
We adopt \citeauthor{MammenTs99}'s
measurement~\citep{MammenTs99} of classification difficulty through the gap
$\Prb(Y = 1 \mid X = x) - 1/2$:
%
for a \emph{noise exponent} $\beta \in (0, \infty)$, we
consider the condition
\begin{equation}
  \label{eqn:mammen-tsybakov}
  \Prb ( | \Prb(Y = 1 \mid X) - 1/2| \leq \epsilon )
  = O(\epsilon^\beta)  .
  \tag{$\textsc{M}_\beta$}
\end{equation}
As $\beta \uparrow \infty$ examples are less likely to have small margin,
and the limit $\beta=\infty$ gives the hard margin $|\P(Y = 1 \mid X) -
\half| \ge \epsilon$ for  small $\epsilon$.
%
Under the independent decomposition
Assumption~\ref{assumption:x-decomposition}, the noise
condition~\eqref{eqn:mammen-tsybakov} solely depends on the covariate's
projection onto the signal $Z$, motivating the following assumption on $Z$.
\begin{assumption}
  \label{assumption:z-density}
  For a given $\beta > 0$, the signal variable
  $Z$ is \emph{$(\beta, c_Z)$-regular}, meaning
  that the absolute value $|Z|$ has density $p(z)$ on $(0, \infty)$, no
  point mass at $0$, and satisfies
  \begin{align*}
    \sup_{z \in (0, \infty)} z^{1 - \beta} p(z) < \infty,
    \qquad \lim_{z \to 0}z^{1 - \beta} p(z) = c_Z \in (0, \infty).
  \end{align*}
\end{assumption}
\noindent
Concretely, when the features $X$ are isotropic Gaussian, $Z
\sim \normal(0, 1)$ and $\beta=1$.

For a link $\sigma$ differentiable at $0$ with
$\sigma'(0) > 0$, we have $\Prb_\sigma(Y = 1 \mid Z = z) = \sigma(t\opt z) =
\half + \sigma'(0) t\opt z + o(t\opt z)$, so the Mammen-Tsybakov noise
condition~\eqref{eqn:mammen-tsybakov} is equivalent to
\begin{align*}
  \Prb \left( \left|t^\star Z \right| \leq \epsilon \right) = O(\epsilon^\beta).
\end{align*}
Thus,
under Assumption~\ref{assumption:z-density},
condition~\eqref{eqn:mammen-tsybakov} holds, as
by dominated convergence%% we have
\begin{align*}
  \Prb\left(|t\opt Z| \le \epsilon\right)
  = \int_0^{\epsilon/t\opt} p(z) dz
  = \int_0^{\epsilon/t\opt} c_Z(1 + o_\epsilon(1)) z^{\beta - 1} dz
  = \frac{c_Z}{\beta} \epsilon^\beta \cdot (1 + o_\epsilon(1)).
\end{align*}
%%   \Prb \left( \left|t^\star Z \right| \leq \epsilon \right) = \int_0^{\frac{\epsilon}{t^\star}} p(z) dz \leq \sup_{z \in (0, \infty)} z^{1 - \beta} p(z) \cdot \int_0^{\frac{\epsilon}{t^\star}} z^{\beta - 1} dz \leq \frac{\sup_{z \in (0, \infty)} z^{1 - \beta} p(z)}{{t^\star}^\beta \beta} \cdot \epsilon^\beta.
%% \end{align*}
We provide extensions of
Corollaries~\ref{cor:lowd-maximum-likelihood-Gaussian}
and~\ref{cor:lowd-majority-vote-m} in the more general cases the noise
exponent $\beta$ allows.
%
The maximum likelihood estimator retains its convergence guarantees in this
setting, and we can generalize the final claim of
Corollary~\ref{cor:lowd-maximum-likelihood-Gaussian} (see
Appendix~\ref{proof:lowd-maximum-likelihood} for a proof):

%% We investigate the non-aggregated MLE estimator and the majority-vote
%% estimator under general distribution assumptions with noise exponent
%% $\beta$.

\begin{proposition}
  \label{prop:lowd-maximum-likelihood}
  Let Assumptions~\ref{assumption:x-decomposition} and
  \ref{assumption:z-density} hold for some $\beta  > 0$ and
  $t\opt = \ltwos{\theta\opt}$. Let $L\lr(\theta) = \Ep[
    \ell_\theta\lr(Y \mid X)]$ be the population loss.
  %
  Then maximum likelihood estimator~\eqref{eq:maximum-likelihood} satisfies
  \begin{equation*}
    \sqrt{n} \left(\what{\theta}\lr_{n,m} - \theta\opt\right)
    \cd \normal(0, m^{-1}\nabla^2 L\lr(\theta\opt)^{-1})  . %\label{eq:mle-normality-theta}
  \end{equation*}
  Moreover,
  \begin{equation*}
    \sqrt{n}\left(\what{u}\lr_{n,m} - u\opt\right)
    \cd \normal\left(0, m^{-1} \cdot (t\opt)^{-2} \proj_{u\opt}^\perp
    \nabla^2 \poploss\lr(\theta\opt)^{-1} \proj_{u\opt}^\perp\right), %\label{eq:mle-normality-u}
  \end{equation*}
  and there exist $C(t)$ and constant $b > 0$ such that
  $C(t) = \frac{b}{ t^{2 - \beta}}(1 + o_t(1))$ and 
  \begin{equation*}
    \sqrt{n} \left(\what{u}\lr_{n,m} - u\opt\right)
    \cd \normal\left(0, m^{-1} \cdot C(t\opt) \prn{\proj_{u\opt}^\perp \Sigma \proj_{u\opt}^\perp}^\dagger \right).  %\label{eq:mle-normality}
  \end{equation*}
\end{proposition}

For majority-vote aggregation, we can generalize
Corollary~\ref{cor:lowd-majority-vote-m} to have
$t_m \asymp \sqrt{m}$, but now
the noise exponent $\beta$ determines the convergence rate
(see Appendix~\ref{proof:lowd-majority-vote-m-noise} for a proof):

\begin{proposition} \label{prop:lowd-majority-vote-m-noise}
  Let Assumptions~\ref{assumption:x-decomposition} and
  \ref{assumption:z-density} hold for some $\beta \in (0, \infty)$
  and $t\opt =
  \ltwos{\theta\opt}$.
  %
  Let $h = h_m$ be the
  function~\eqref{eqn:h-centering} with signal covariate $Z$ defined in
  Assumption~\ref{assumption:x-decomposition}.
  %
  Then there are constants $a, b > 0$, depending only on $\beta$ and $c_Z$,
  such that the following hold: there is a unique $t_m \geq t_1 = t^\star$
  solving $h(t_m) = 0$ and for this $t_m$ we have both
	\begin{equation*}
		\what{\theta}\mv_{n,m} \cp t_m u\opt
		~~ \mbox{and} ~~
		\lim_{m \to \infty} \frac{t_m}{t\opt \sqrt{m}} = a.
	\end{equation*}
	Moreover,
	there exists a function
	$C_m(t) = \frac{b}{(t \sqrt m)^{2 - \beta}}(1 + o_m(1))$
	as $m \to \infty$ such that
	\begin{align*}
		\sqrt{n} \left(\what{u}\mv_{n,m} - u^\star \right)
		&\cd \normal\left(0, C_m(t\opt) \prn{\proj_{u\opt} \Sigma \proj_{u\opt}}^\dagger\right).
	\end{align*}
\end{proposition}

Paralleling the discussion in Section~\ref{sec:gauss}, we may compare the
performance of the MLE $\what{\theta}\lr_{n,m}$ and
the majority-vote estimator $\what{\theta}\mv_{n,m}$.
%
When the classification problem is hard---meaning that
$\beta$ in Condition~\eqref{eqn:mammen-tsybakov} is near 0%% so that
%%that classifying most examples is nearly random chance
---the aggregation in the majority vote estimator still allows
convergence (nearly) as quickly as the non-aggregated estimator; the
problem is so noisy that data ``cleaning'' by aggregation is helpful.
For easier problems, where $\beta \gg 0$, the gap between them
grows; this is sensible, as aggregation is likely to
force a dataset to be separable, thus making fitting methods
unstable (a minimizer may even fail to exist). 

%% Then we see that $\what{\theta}\mv_{n,m}$ is always not calibrated where
%% $\what{\theta}\mv_{n,m} \to \Theta(\sqrt{m} \theta\opt)$ regardless of the
%% noise exponent, while the MLE estimator is calibrated. In terms of
%% classification, the majority-vote estimator also exhibits worse convergence
%% rates $\frac{O(1)}{\sqrt{n m^{1 - \frac{1}{2} \beta}}}$ in
%% Theorem~\ref{thm:lowd-majority-vote-m-noise} comparing to
%% $\frac{O(1)}{\sqrt{nm}}$ in Theorem~\ref{thm:lowd-maximum-likelihood}. We
%% recover our results in Sec.~\ref{sec:non-agg-gauss} by taking $\beta = 1$,
%% and we also see that for noisier data with smaller $\beta$, the
%% majority-vote estimator is more effective.

%% \cc{Perhaps discussion about $\beta > 2$? If we don't reviewers might ask}


